\documentclass{article}
\usepackage{graphicx, amsfonts, amsmath, amssymb, amsthm, float, cancel, braket} % Required for inserting images

\title{Struttura della Materia}
\author{Alessandro Ballini}
\date{March 2026}

\begin{document}

\maketitle
\tableofcontents
\newpage

\section*{Conversioni}
Vediamo delle conversioni per alcune grandezze fisiche che risultano essere comode rispetto alle scale caratteristiche dei sistemi in esame. Riportiamo, prima di tutto, le grandezze fondamentali:
\begin{align}
    &\text{Massa elettrone $e$} \qquad &m_e = 9.109 \cdot 10^{-31} \, kg \notag \\
    &\text{Massa protone $p$} \qquad &m_p = 1.673 \cdot 10^{-27} \, kg \notag \\
    &\text{Massa neutrone $n$} \qquad &m_n = 1.675 \cdot 10^{-27} \, kg \notag \\
    \notag \\
    &\text{Carica fondamentale} \qquad &e = 1.602 \cdot 10^{-19} \, C\cdot m \notag \\
    \notag \\
    &\text{Costante di Planck $h$} \qquad &h = 6.62 \cdot 10^{-34} \, \frac{J}{s} \notag \\
    &\text{Costante di Planck ridotta $\hbar$} \qquad &\hbar = 1.05 \cdot 10^{-34} \frac{J}{s} \notag \\
    \notag \\
    &\text{Costante dielettrica nel vuoto $\varepsilon_0$} \qquad &\varepsilon_0 = 8.85 \cdot 10^{-12} \, \frac{C^2}{N\cdot m^2} \notag \\
    &\text{Permeabilità magnetica nel vuoto $\mu_0$} \qquad &\mu_0 = 4\pi \cdot 10^{-7} \frac{H}{m} \notag \\
    \notag \\
    &\text{Velocità della luce $c$} \qquad &c = 3 \cdot 10^8 \, \frac{m}{s} \notag \\
    \notag \\
    &\text{Angstrong $\mathring{A}$} \qquad &1 \, \mathring{A} = 10^{-10} \, m \notag \\
    &\text{Fermi $fm$} \qquad &1 \, fm = 10^{-15} \, m \notag \\
    \notag \\
    &\text{Spettro del visibile} \qquad &\lambda = (400 \div 800) \cdot 10^{-9} \, m \notag
\end{align}
L'energia $\mathcal{E}$, di solito espressa in Joule $J$, può essere anche espressa in electron-volts $eV$ o in funzione dell'inverso di una lunghezza, come ad esempio $\mathring{A}^{-1}$ o $cm^{-1}$. In particolare $1 \, eV$ è l'energia che ha un elettrone sottoposto ad una differenza di potenziale di $1 \, V$ (in modulo), ovvero
\begin{align}
    1 \, eV = 1.602 \cdot 10^{-19} J \notag
\end{align}
Usando la relazione di Einstein $E = \hbar\omega = h\nu = \frac{hc}{\lambda}$ possiamo esprimere l'energia funzione dell'inverso di una lunghezza d'onda. In particolare se esprimiamo la lunghezza d'onda in Angstrom si ha
\begin{align}
    E(eV) &= \frac{1}{e} \cdot \frac{h(J/s) \cdot c(\mathring{A}/s)}{\lambda(\mathring{A})} \notag \\
    &= \frac{1}{1.602 \cdot 10^{-19} C\cdot m}\frac{6.62 \cdot 10^{-34} \, J/s \cdot 3 \cdot 10^{18} \, \mathring{A}/s}{\lambda(\mathring{A})} \notag \\
    &= \frac{12.4 \cdot 10^3 \frac{\mathring{A} \cdot J}{C \cdot m \cdot s^2}}{\lambda(\mathring{A})} \notag
\end{align}
ovvero
\begin{gather}
    E(eV) = \frac{12.4 \cdot 10^3}{E(\mathring{A}^{-1})} \qquad E(\mathring{A}^{-1}) = \frac{12.4 \cdot 10^3}{E(eV)} \notag
\end{gather}
Usando questa relazione si ottiene la larghezza dello spettro del visibile espressa in $eV$, ovvero
\begin{gather}
    \text{Spettro del visibile} \qquad (1.55 \div 3.10) \, eV \notag
\end{gather}
dove $1.55 \, eV$ è il rosso e $3.10 \, eV$ è il blu (dato che l'energia è inversamente proporzionale alla lunghezza d'onda).

\newpage
\section*{Sistema di unità atomiche (u.a.)}
Un sistema molto comodo per svolgere calcoli riguardanti sistemi atomici e molecolari è il sistema di unità atomiche (u.a.). Tale sistema pone arbitrariamente
\begin{align}
    1 &= \hbar \notag \\
    1 &= 4\pi\varepsilon_0 \notag \\
    1 &= e \notag \\
    1 &= m_e \notag
\end{align}
segue che il raggio di Bohr $a_0 = 1 \, u.a.$. Possiamo ricavare le conversione inverse, ovvero vedere la relazione tra unità atomiche e unità fisiche:
\begin{align}
    &u.a.(l) \quad &= 0.529 \, \mathring{A} \notag \\
    &u.a.(m) \quad &= 9.1 \cdot 10^{-31} \, kg \notag \\
    &u.a.(v) \quad &= 2.2 \cdot 10^6 \, \frac{m}{s} \notag \\
    &u.a.(t) \quad &= \frac{0.529 \, \mathring{A}}{2.2 \cdot 10^6 \, m/s} = 2.4 \cdot 10^{-17} \, s \notag \\
    &u.a.(\mathcal{E}) \quad &= -2\mathcal{E}^H_{1s} = 27.21 \, eV \notag \\
    &u.a.(B) \quad &= 2.35 \cdot 10^5 \, T \notag \\
    &u.a.(E) \quad &= 5.13 \cdot 10^{11} \, \frac{V}{m} \notag \\
    &u.a.(d) \quad &= 8.46 \cdot 10^{-30} C \cdot m \notag
\end{align}
Inoltre si ottiene dal fatto che la costante di struttura affine $\alpha = 1/137$ che $c = 137 \, u.a.$.

\newpage
\section{Atomo idrogenoide}
\subsection{Approccio semiclassico}
Consideriamo un atomo idrogenoide con numero atomico $Z$. L'elettrone, a distanza $r$ dal nucleo, è soggetto alla forza di Coulomb e alla forza centripeta. Affinché l'elettrone resti in orbita la risultante delle forze deve essere identicamente nulla, ovvero deve valere
\begin{gather}
    \frac{m_ev^2}{r} = \frac{Ze^2}{4\pi\varepsilon_0r^2} \notag
\end{gather}
dunque si ottiene
\begin{gather}
    m_ev^2 = \frac{Ze^2}{4\pi\varepsilon_0r} \notag
\end{gather}
Questo ci permette di ricavare un'espressione per l'energia cinetica $T$ dell'elettrone
\begin{gather}
    T = \frac{1}{2}m_ev^2 = \frac{1}{2}\frac{m_eZe^2}{4\pi\varepsilon_0r} \notag
\end{gather}
Inoltre, l'energia potenziale a cui è sottoposto l'elettrone è nota, infatti
\begin{gather}
    V = -\frac{Ze^2}{4\pi\varepsilon_0r} \notag
\end{gather}
notiamo che $V = -2T$. Esprimendo l'energia totale dell'elettrone si ha
\begin{gather}
    \mathcal{E} = T + V = -T = -\frac{1}{2}\frac{m_eZe^2}{4\pi\varepsilon_0r} \notag
\end{gather}
Questo modello non può essere corretto, infatti l'elettrone irradia energia durante il moto e dovrebbe cadere, in un tempo dell'ordine di $10^{-11} \, s$, nel nucleo. Tuttavia questo contraddice la realtà, infatti gli atomi sono stabili e l'elettrone non cade nel nucleo. Niels Bohr ipotizzò che il momento angolare dell'elettrone non potesse assumere valori arbitrari ma fosse vincolato ad essere un multiplo intero di $\hbar$ per spiegare le osservazioni sperimentali. Dunque viene aggiunta un'altra condizione, in particolare
\begin{gather}
    \begin{cases}
        m_ev^2_n &= \frac{Ze^2}{4\pi\varepsilon_or_n} \\
        mv_nr_n &= n\hbar
    \end{cases} \notag
\end{gather}
questo introduce dei vincoli anche su $v$ e $r$, dunque anch'essi ora dipendono da $n$. Risolvendo il sistema si ottiene
\begin{gather}
    \begin{cases}
        \frac{n^2\hbar^2}{m_er^2_n} &= \frac{Ze^2}{4\pi\varepsilon_0r_n} \\
        v_n &= \frac{n\hbar}{m_er_n}
    \end{cases} \longrightarrow
    \begin{cases}
        r_n &= \frac{4\pi\varepsilon_0}{Ze^2}\frac{n^2\hbar^2}{m_e} \\
        v_n &= \frac{Ze^2}{4\pi\varepsilon_0}\frac{1}{n\hbar}
    \end{cases} \notag
\end{gather}
Usando la relazione prima ricavata, $\mathcal{E} = -T$,  si ottengono i corretti livelli energetici dell'elettrone, i raggi delle orbite per le quali si ha moto stabile (niente irraggiamento) e la velocità dell'elettrone in tale orbita. In particolare
\begin{align}
    \mathcal{E}^H_n = -\frac{1}{2}m_e\left(\frac{Ze^2}{4\pi\varepsilon_0}\right)^2\frac{1}{n^2\hbar^2} \notag
\end{align}
inoltre $r^H_{n = 1} = a_0$ e $v^H_{n = 1} = 2.2 \cdot 10^6 m/s$.

\subsection{Correzioni relativistiche}
L'energia relativistica dell'elettrone in un atomo idrogenoide è
\begin{align}
    E_{rel} &= \sqrt{p^2c^2 + m^2_ec^4} \notag \\
    &= m_ec^2\sqrt{1 + \frac{p^2}{m^2_ec^2}} \notag
\end{align}
dunque usando l'approssimazione $\sqrt{1 + \varepsilon} \sim 1 + \frac{1}{2}\varepsilon - \frac{1}{8}\varepsilon^2$ si ottiene
\begin{align}
    E_{rel} &\sim m_ec^2\left(1 + \frac{p^2}{2m^2_ec^2} - \frac{p^4}{8m^4_ec^4}\right) \notag \\
    &= m_ec^2 + \frac{p^2}{2m_e} - \frac{p^4}{8m^3_ec^2} \notag
\end{align}
pertanto, trascurando il termine costante $m_ec^2$ che è l'energia a riposo dell'elettrone libero, si ha che, in prima approssimazione, la correzione all'hamiltoniana è data dal termine
\begin{align}
    H'_{k} &= -\frac{p^4}{8m^3_2c^2} \notag \\
    &= -\frac{1}{2m_ec^2}T^2 \notag \\
    &= -\frac{\alpha^2}{2}T^2 \qquad (u.a.) \notag
\end{align}
dove $T$ è l'operatore energia cinetica. Possiamo allora applicare i metodi della teoria delle perturbazioni sull'hamiltoniano perturbato in termini di $\alpha$. Notiamo che, nonostante il fatto che l'atomo idrogenoide presenti degenerazione, possiamo applicare il metodo delle perturbazioni indipendenti dal tempo nel caso non degenere, infatti la perturbazione commuta con l'hamiltoniano imperturbato e quindi è già diagonale nella base iniziale (non "mischia" gli autostati del sottospazio degenere). Allora si ha
\begin{align}
    \Delta E^k_{nlm} &= \braket{\psi^{(0)}_{nlm}|-\frac{\alpha^2}{2}T^2|\psi^{(0)}_{nlm}} \notag
\end{align}
In particolare $T = H + \frac{Z}{r}$ per l'atomo idrogenoide, quindi
\begin{align}
    T^2 = H^2 + \frac{Z^2}{r^2} + 2H\frac{Z}{r} \notag
\end{align}
e sostituendo si ottiene
\begin{align}
    \Delta E^k_{nlm} &= -\frac{\alpha^2}{2}(E^2_n + Z^2\braket{1/r^2}_{nlm} + 2E_nZ\braket{1/r}_{nlm}) \notag
\end{align}
dove
\begin{align}
    \braket{1/r}_{nlm} &= \int_0^\infty rR(r)_{nl} \, dr \notag \\
    \braket{1/r^2}_{nlm} &= \int_0^\infty R(r)_{nl} \, dr \notag
\end{align}
dunque in unità atomiche si ha
\begin{align}
    \Delta E^k_{nlm} &= -\frac{\alpha^2}{2}\left(\frac{Z^4}{4n^4} - \frac{Z^5}{2n^4} + \frac{Z^4}{n^3\left(l + \frac{1}{2}\right)}\right) \notag \\
    &= -E_n\left(\frac{Z\alpha}{n}\right)^2\left(\frac{3}{4} - \frac{n}{l - \frac{1}{2}}\right) \leq 0 \notag
\end{align} \\
La seconda correzione è dovuta al moto del nucleo del sistema di riferimento dell'elettrone, infatti si osserva una densità di corrente non nulla dovuta al moto relativo del nucleo e dunque un campo magnetico. In particolare si ha 
\begin{align}
    \mathbf{j} = -Ze\mathbf{v} \notag
\end{align}
Il campo magnetico osservato è dato da
\begin{align}
    \mathbf{B} = \frac{\mu_0}{4\pi}\frac{\mathbf{j} \times \mathbf{u}_r}{r^2} = -\frac{\mu_0}{4\pi}Ze\frac{\mathbf{v} \times \mathbf{u}_r}{r^2} \notag
\end{align}
ponendo
\begin{align}
    \mathbf{E} = \frac{Ze^2}{4\pi\varepsilon_0r^2}\mathbf{u}_r \notag
\end{align}
possiamo scrivere
\begin{align}
    \mathbf{B} = -\frac{1}{c^2}\mathbf{v} \times \mathbf{E} \notag
\end{align}
inoltre, per simmetria radiale, si ha
\begin{align}
    \mathbf{E} = \frac{1}{r}\frac{dV}{dr}\mathbf{u}_r \notag
\end{align}
pertanto la correzione all'hamiltoniana diventa
\begin{align}
    H'_{so} &= -\mathbf{\mu}_S \cdot \mathbf{B} \notag \\
    &= -\frac{g_s\mu_B}{\hbar}\mathbf{S} \cdot \mathbf{B} \notag \\
    &= \frac{\alpha^2}{2}\frac{1}{r}\frac{dV}{dr}\mathbf{L} \cdot \mathbf{S} \notag
\end{align}
in unità atomiche. Allora 
\begin{align}
    \Delta E^{so}_{nlm} &= \frac{\alpha^2}{2}\braket{\psi^{(0)}_{nlm}|\frac{1}{r}\frac{dV}{dr}\mathbf{L} \cdot \mathbf{S}|\psi^{(0)}_{nlm}} \notag \\
    &= \frac{\alpha^2\hbar^2}{4}\braket{Z/r^3}_{nlm} \cdot (j(j + 1) - l(l + 1) - s(s + 1)) \notag
\end{align} \\
Infine abbiamo la correzione per lo stato $1s$ dovuta al fatto che la probabilità di trovare lo stato $l = 0$ nel nucleo è non nulla, ovvero $|R_{10}(0)|^2 \neq 0$. Sia $U(r)$ il potenziale di interzione tra il nucleo e l'elettrone dello stato $1s$ nel sistema di riferimento dell'elettrone. Per il principio di indeterminazione l'orbita dell'elettrone non è ben identificata, si ha 
\begin{align}
    \delta r \sim \frac{\hbar}{\delta p} \sim \alpha \notag
\end{align}
allora 
\begin{align}
    U(r + \delta r) &\sim U(r) + \delta r \nabla U + \sum_{i, j}\delta x_i \delta x_j \frac{\partial^2 U}{\partial x_i \partial x_j} \notag
\end{align}
ovvero 
\begin{align}
    dU \sim \frac{\alpha^2}{6}\nabla^2 U \notag
\end{align}
dove si è usato $\delta x_i \delta x_j = \delta_{ij}\frac{\delta r^2}{3}$. Inoltre $V(r) = - U(r)$ e quindi, poiché $\nabla^2V = 4\pi\rho(r) \sim 4\pi\delta(r)$ (l'approssimazione è dovuta al fatto che le dimensioni del nucleo sono molto più piccole rispetto alla distanza media dell'elettrone dal nucleo stesso), si ottiene $\nabla^2U \sim -4\pi\delta(r)$, ovvero
\begin{align}
    dU \sim \frac{2}{3}\alpha^2\pi\delta(r) \notag
\end{align}
e quindi
\begin{align}
    \Delta E^{D}_{nlm} &\sim \braket{\psi^{(0)}_{n00}|\frac{2}{3}Z\alpha^2\pi\delta(r)|\psi^{(0)}_{n00}} \notag \\
    &= \frac{2}{3}Z\alpha^2\pi|\psi^{(0)}_{n00}(0)|^2 \notag \\
    &= -E_n\frac{4}{3}\frac{(Z\alpha)^2}{n} \notag
\end{align}
In realtà la correzione giusta, applicando la QCD, sarebbe
\begin{align}
    H'_D &= \frac{Z\alpha^2}{2}\pi\delta(r) \notag \\
    \Delta E^D_{nlm} &= -E_n\frac{(Z\alpha)^2}{n} \notag
\end{align}

\subsubsection{Esercizi}
\textbf{1.} Determinare l'errore che si compie sul calcolo dell'energia dello stato fondamentale dell'atomo di idrogeno se si sceglie come autofunzione dello stato fondamentale la seguente
\begin{align}
    \psi(r, \theta, \varphi) = Ne^{-cr^2} \notag
\end{align}
con $c > 0$. \\

\textbf{Soluzione:} la costante $N$ si determina imponendo la normalizzazione della funzione d'onda su $\mathbf{R}^3$. Per determinare l'energia dello stato fondamentale utilizzando $\psi$ come funzione d'onda procediamo calcolando il valor medio di $H$ rispetto a $\psi$ e cercando per quale $c$ positivo si ha che il valore medio di $H$ è minimo. In particolare imponiamo
\begin{align}
    4\pi N^2\int_0^\infty r^2e^{-2cr^2} \, dr = 1 \notag
\end{align}
allora otteniamo 
\begin{align}
    N =   \notag
\end{align}
A questo punto $\braket{H}_c = \braket{T}_c + \braket{V}_c$, quindi
\begin{align}
    \braket{V}_c &= 4\pi N^2 \int_0^\infty r^2e^{-2cr^2}\frac{1}{r} \, dr \notag \\
    \braket{T}_c &= -2\pi N^2 \int_0^\infty r^2\psi^*(r, \theta, \varphi)\nabla^2\psi(r, \theta, \varphi) \, dr \notag
\end{align} \\ \\
\textbf{2.} Consideriamo il seguente hamiltoniano 
\begin{align}
    H_\mu = -\frac{1}{2\mu}\nabla^2 - \frac{1}{r} \notag
\end{align}
con $\mu > 0$, determinare l'errore sul calcolo dell'energia dello stato $1s$ usando il metodo perturbativo e variazionale lineare. \\

\textbf{Soluzione:} riscriviamo
\begin{align}
    H_\mu = H_1 + H'_\mu \notag
\end{align}
dove $H_1$ è l'hamiltoniano imperturbato dell'atomo di idrogeno, ovvero 
\begin{align}
    H_1 = -\frac{1}{2}\nabla^2 - \frac{1}{r} \notag
\end{align}
mentre $H'_\mu$ è la perturbarzione introdotta con il parametro $\mu$, ovvero 
\begin{align}
    H'_\mu = \left(1 - \frac{1}{\mu}\right)\frac{\nabla^2}{2} \notag
\end{align}

\newpage
\section*{Appendice}
Per ogni $k \in \mathbf{N}$ si ha
\begin{align}
    \int_0^\infty x^{2k + 1}e^{-ax^2} \, dx &= \frac{k!}{2a^{k + 1}} \notag \\
    \int_{-\infty}^\infty x^{2k}e^{-ax^2} \, dx &= \frac{(2k - 1)!!}{(2a)^k}\sqrt{\frac{\pi}{a}} \notag \\
    \int_0^\infty x^{2k}e^{-ax^2} \, dx &= \frac{(2k - 1)!!}{2(2a)^k}\sqrt{\frac{\pi}{a}} \notag
\end{align}



\end{document}